\documentclass[11pt,a4paper]{article}
\usepackage[margin=1in]{geometry}
\usepackage{amsmath,booktabs,hyperref,float}

\title{Iteration 025: Cosine Loss Schedule}
\author{Autoresearch Loop}
\date{April 8, 2026}

\begin{document}
\maketitle

\begin{abstract}
A cosine-annealed loss schedule from 80\% MSE to 20\% MSE over 150 epochs,
combined with InstanceNorm, achieves $r=0.376$, not improving over the
fixed 50/50 loss ($r=0.378$). InstanceNorm degrades performance and masks
any benefit from loss scheduling.
\end{abstract}

\section{Hypothesis}
Early training benefits from MSE (stable gradients for large initial errors),
while late training should emphasize correlation loss (optimizes the actual
metric). A cosine schedule from $\alpha=0.8$ (MSE-heavy) to $\alpha=0.2$
(corr-heavy) gives the model stable convergence then metric-optimal refinement.

\section{Method}
$\alpha(t) = 0.2 + 0.5(0.8-0.2)(1+\cos(\pi t/T))$, decaying from 0.8 to 0.2.
Architecture: InstanceNorm + FIR (7 taps). Corr validation for early stopping.

\section{Results}
\begin{table}[H]
\centering
\begin{tabular}{lrrr}
\toprule
Model & Mean $r$ & Std $r$ & SNR (dB) \\
\midrule
Combined loss (iter017) & 0.378 & 0.078 & 0.61 \\
Cosine schedule + InstanceNorm & 0.376 & 0.077 & 0.63 \\
\bottomrule
\end{tabular}
\end{table}

\section{Analysis}
The cosine loss schedule shows the highest SNR of any model tested (0.63\,dB),
suggesting the loss scheduling helps with amplitude prediction. However, the
lower $r$ indicates InstanceNorm is the dominant factor. The loss scheduling
idea may have merit without InstanceNorm --- worth revisiting.

\section{Next}
Clean variants: mixup, pure corr loss, and high channel dropout --- all without InstanceNorm.

\end{document}
